%%%%%%%%%%%%%%%%
% ABSTRACT
%%%%%%%%%%%%%%%%

\section*{Abstract}

Architectural drift is one of the main causes of structural software deterioration in SMEs, because design decisions progressively diverge from the actual implementation when no continuous verification mechanism exists. This proposal aims to develop and experimentally evaluate an automated mechanism for detecting architectural drift based on static code analysis, dependency graph construction, and the execution of \textit{Fitness Functions}, in order to support technical governance without relying on external architectural documentation. The work will be limited to a single technological ecosystem and will evaluate three specific conformance dimensions: layer integrity, absence of dependency cycles, and control of excessive outgoing dependency concentration through explicit thresholds. Based on these rules, the mechanism will generate a structural representation of the system, calculate a Conformance Rate (CR), and classify the architectural state as acceptable, under observation, or non-conforming. The experimental demonstration will compare a baseline version of the system with altered versions containing controlled drifts. The expected outcomes include a functional prototype, a documented rule set, and evidence of the usefulness of the approach to support maintenance and refactoring decisions in agile development environments. In this way, the project seeks to contribute a practical and novel governance mechanism for SMEs by combining code-only evidence, dependency graphs, executable architectural rules, and explicit conformance thresholds.


